\section{Time, space, and bit complexity}
\label{sec:complexity}

\subsection{Arithmetic complexity of differential CMP}

Let \(K=\max_iK_i\) and \(k=P+K\).  Suppose every evaluated Jacobian block is
an explicitly supplied matrix over a field and has dimensions \(O(k)\).
Classical dense elimination computes a null-space basis, multiplies the
surviving blocks, and row-reduces the result using \(O(k^3)\) field operations
per bag.  Therefore
\begin{equation}
  T_{\mathrm{arith}}(N,P,K)=O\bigl(N(P+K)^3\bigr).
  \label{eq:arithmetic-complexity}
\end{equation}
This is a valid bound for \cref{thm:differential-cmp-correct}.

If \(N,P,K\le n\), then
\begin{equation}
  N(P+K)^3
  \le n(2n)^3
  =8n^4.
  \label{eq:uniform-arithmetic-inequality}
\end{equation}
Thus the number of field operations is uniformly \(O(n^4)\) once the matrices
are present.  For fixed \(P,K\), \cref{eq:arithmetic-complexity} is linear in
\(N\).

\subsection{Space}

A streaming implementation that processes and discards leaf matrices uses
\begin{equation}
  O\bigl(V+E+(P+K)^2\bigr)
  \label{eq:streaming-space}
\end{equation}
field elements in addition to constant bookkeeping per bag, where \(V+E\)
denotes the explicitly stored constraint incidence structure.  Retaining all
messages uses
\begin{equation}
  O\bigl(V+E+N(P+K)^2\bigr)
\end{equation}
field elements.

\subsection{Why arithmetic counting is not a P theorem}

Membership in \(P\) is a bit-complexity statement.  The following data must be
specified and bounded:

\begin{enumerate}
  \item the representation of every coefficient;
  \item the bit length of intermediate matrix entries;
  \item the cost of exact zero testing and rank decisions;
  \item the representation of any algebraic real at which a Jacobian is
  evaluated; and
  \item the size and construction cost of the zeroth-order message needed for
  global feasibility.
\end{enumerate}

For an explicitly given rational matrix, fraction-free elimination controls
coefficient growth and yields polynomial bit complexity in the matrix
dimensions and input bit length~\cite{Bareiss1968}.  This fact does not close
the present proof: a feasible point can have irrational coordinates and is not
part of the input, while finding or deciding the existence of such a point is
the original nonlinear problem.

Similarly, treating symbolic entries as elements of a multivariate quotient
ring changes the operation model.  Rank, kernels, and equality testing over
that ring cannot be charged as unit-cost field operations.  Their
representations can require polynomial reduction, localization, or ideal
membership computations.

\begin{theorem}[Conditional bit-complexity implication]
\label{thm:conditional-bit}
Assume there is an exact CMP representation satisfying
\cref{prob:polynomial-semantic-message}, with total bit running time
\begin{equation}
  N\,\poly(P+K,L)
  \label{eq:desired-bit-bound}
\end{equation}
and a polynomial-time root nonemptiness test.  Then
\(\ExactNNTP\in P\).
\end{theorem}

\begin{proof}
The instance dimensions and coefficient lengths are bounded by its explicit
binary encoding length.  Under the assumed bound, the stated CMP procedure is
a deterministic polynomial-time Turing decider whose acceptance semantics is
exactly \(\Feas(I)\).
\end{proof}

\begin{remark}[Current status]
The hypothesis of \cref{thm:conditional-bit} is \openstatus.  Equation
\cref{eq:uniform-arithmetic-inequality} proves only the arithmetic inequality;
it does not construct a global decider.
\end{remark}
